% Appendix: proofs and negative results for Section "Beyond dictatorship" (sec:smd).
% Requires \appendix before inclusion.
\section{The SMD mechanism: proofs and variants}
\label{app:smd}

Throughout, $\tau=(1+\varepsilon)/2$, agent~1 has type $x$ and agent~2 type $y$. In the coin-flip branch
the dictator takes $(s,0)$ if her type is at least $\tfrac12$ and $(0,s)$ otherwise. If the dictator
of type $d$ takes $(s,0)$ and the other agent has type $o$, welfare is
$ds+o(1-s)+(1-o)+\varepsilon(1-s)=1+\varepsilon+s(d-o-\varepsilon)$, and symmetrically for $(0,s)$.

\paragraph{Proof of Proposition~\ref{prop:D}.} Consider agent~1 (agent~2 is symmetric). If $1-\tau<y<\tau$, no report
of agent~1 triggers a split, the outcome is the coin flip, and a fixed menu is strategy-proof
(Proposition~\ref{prop:menu}). Let $y\le1-\tau$ (the case $y\ge\tau$ is the mirror image). Reporting
$r\ge\tau$ triggers the split and gives agent~1 item~1, worth $x$. Any report $r<\tau$ gives the coin
flip, in which she can still take $(s,0)$ or $(0,s)$ as dictator, and as the other agent she receives
$(1,1-s)$. For $x\ge\tfrac12$ the best coin-flip utility is
\[
F(x)=\tfrac12xs+\tfrac12\big[x+(1-x)(1-s)+\varepsilon(1-s)\big]=xs+\tfrac12(1-s)(1+\varepsilon),
\]
so $F(x)-x=(1-s)(\tau-x)$: the split is better exactly when $x\ge\tau$, which is when truthful
reporting triggers it. For $x<\tfrac12$ the split gives item~1, worth $x<\tfrac12$, while truthful
reporting gives at least $\tfrac12[x+(1-x)+\varepsilon(1-s)]\ge\tfrac12$. Within the coin flip the
choice is from a fixed menu. Hence no report beats the truth. A brute-force check over a grid of
true types, reports and opponents, for several $s$ and $\varepsilon$, finds no profitable deviation. \qed

\paragraph{Proof of Proposition~\ref{prop:E}.} Take $x\ge y$. In the split region $x-y\ge\varepsilon$, so the first best
is the split and the ratio is $1$. Outside it:
\begin{itemize}[nosep]
\item \emph{Same side} ($x,y\ge\tfrac12$, or both below): both dictators take $s$ of the same item, so
  $W=1+\varepsilon-\varepsilon s$ while $\mathrm{OPT}\le1+\max\{\varepsilon,\tfrac12\}$ (the gap is at most
  $\tfrac12$); the ratio is at least $(1+\varepsilon-\varepsilon s)/(1+\max\{\varepsilon,\tfrac12\})$, with
  equality at identical types when $\varepsilon\ge\tfrac12$ and at $(1,\tfrac12)$ when $\varepsilon<\tfrac12$.
\item \emph{Opposite sides, not both strong}: $W=1+\varepsilon+s(\Delta-\varepsilon)$ with $\Delta=x-y$. For
  $\Delta\le\varepsilon$ the ratio is at least the same-side value. For $\Delta>\varepsilon$ it is
  $[1+\varepsilon+s(\Delta-\varepsilon)]/(1+\Delta)$, decreasing in $\Delta$ (the derivative has the sign of
  $-(1-s)(1+\varepsilon)$), and the largest gap outside the split region is $\Delta=\tau$, giving
  $[2(1+\varepsilon)+(1-\varepsilon)s]/(3+\varepsilon)$.
\end{itemize}
The first bound decreases and the second increases in $s$. Equalising them gives, for
$\varepsilon\ge\tfrac12$, $s^\star=(1-\varepsilon^2)/(1+3\varepsilon)$ and ratio $(1+\varepsilon)^2/(1+3\varepsilon)$; for
$\varepsilon\le\tfrac12$, $s^\star=\varepsilon(1+\varepsilon)/(\tfrac32+\tfrac32\varepsilon+\varepsilon^2)$ and ratio
$2(1+\varepsilon)^2/(3+3\varepsilon+2\varepsilon^2)$. The grid minimum of the ratio agrees with these values to
four decimals for every $\varepsilon$ we tested. \qed

\paragraph{Why the rule is rigid.}
\begin{itemize}[nosep]
\item \emph{Probabilistic split for moderately opposed agents.} Suppose agents on opposite sides
  who are not both strong are split with probability $p>0$. Against a strongly opposed opponent, a
  type $x=\tfrac12$ gets $\tfrac p2+(1-p)\tfrac12[1+\varepsilon(1-s)]$ by reporting truthfully, but
  $\tfrac12[1+\varepsilon(1-s)]$ by claiming to prefer item~2 (which removes the split); the deviation
  gains $\tfrac12p\,\varepsilon(1-s)>0$, and by continuity so do types slightly above $\tfrac12$. Moderate
  types must therefore be spared the split entirely.
\item \emph{Larger threshold.} With $\tau'>(1+\varepsilon)/2$, types in $[(1+\varepsilon)/2,\tau')$ strictly
  prefer a split they cannot trigger and misreport upward; on a grid this also lowers the best
  achievable ratio.
\item \emph{Richer fallbacks.} Using nested LIP as the coin-flip menu violates incentive
  compatibility at the threshold (its fallback utility at $\tau$ exceeds $\tau$; gains up to $0.09$).
  Optimising the fallback menu by an LP, given the split region, and evaluating the result exactly
  on the continuum improves on SMD only for $\varepsilon\lesssim0.45$ ($0.857$ vs $0.806$ at
  $\varepsilon=\tfrac14$, $0.885$ vs $0.867$ at $0.4$); at $\varepsilon=0.6$ the LP value equals
  $\rho_{\mathrm{SMD}}$, so SMD is optimal within this class there.
\end{itemize}
